\answerAnchor{MATH1-CALC-0062}
\begin{solutionBox}
\textbf{A}\quad
\[
\overrightarrow{PQ}=(3,-4,2),\qquad |PQ|=\sqrt{29},\qquad
\vec e=\frac1{\sqrt{29}}(3,-4,2).
\]
核验：\(|\vec e|=1\)。

\textbf{B}\quad
\[
\lambda^2+4+(\lambda-1)^2=9
\Longleftrightarrow \lambda^2-\lambda-2=0,
\]
故 \(\lambda=2\) 或 \(\lambda=-1\)。两值代回所得向量均非零且模为 \(3\)。

\textbf{C}\quad
\(\overrightarrow{AB}=(-2,1,2)\)，模为 \(3\)，故单位向量为
\[
\left(-\frac23,\frac13,\frac23\right).
\]

\textbf{M}\quad
\[
\overrightarrow{PQ}=(-k,1,2),\qquad k^2+5=9,
\]
故 \(k=\pm2\)。其同向单位向量第一分量为 \(-k/3\)，要求其为正，
所以 \(k=-2\)。代回所得单位向量为 \((2,1,2)/3\)。
\end{solutionBox}

\answerAnchor{MATH1-CALC-0063}
\begin{solutionBox}
\textbf{A}\quad
\[
3\vec a+2\vec b=(6,-3,0)+(-2,6,4)=(4,3,4).
\]

\textbf{B}\quad
若 \((2,\lambda,6)=\mu(1,-2,3)\)，第一分量给 \(\mu=2\)，
第三分量相容，第二分量给 \(\lambda=-4\)。

\textbf{C}\quad
以 \(A\) 为起点，共面条件为
\[
\det\begin{pmatrix}1&0&1\\0&2&1\\2&2&m\end{pmatrix}
=2m-6=0,
\]
故 \(m=3\)。此时
\(\overrightarrow{AD}=2\overrightarrow{AB}+\overrightarrow{AC}\)，可交叉核验。

\textbf{M}\quad
设 \(\vec c=\alpha\vec a+\beta\vec b\)，则
\[
(1,t,2)=(\alpha,\alpha+\beta,\beta).
\]
所以 \(\alpha=1,\beta=2,t=3\)，且
\(\vec c=\vec a+2\vec b\)。
\end{solutionBox}

\answerAnchor{MATH1-CALC-0064}
\begin{solutionBox}
\textbf{A}\quad
点积为 \(-1\)，两模均为 \(\sqrt2\)，故
\[
\cos\theta=-\frac12,\qquad \theta=\frac{2\pi}{3}.
\]

\textbf{B}\quad
\(\vec a\cdot\vec b=4,|\vec b|=\sqrt5\)，故
\[
\text{标量投影}=\frac4{\sqrt5},\qquad
\operatorname{proj}_{\vec b}\vec a
=\frac45(2,0,1)=\left(\frac85,0,\frac45\right).
\]

\textbf{C}\quad
\[
(\lambda,1,2)\cdot(1,-2,1)=\lambda=0,
\]
故 \(\lambda=0\)。

\textbf{M}\quad
令平方距离
\[
\phi(t)=(2-t)^2+(1-t)^2+4
=2\left(t-\frac32\right)^2+\frac92.
\]
故 \(t=3/2\) 时最小，最小值为
\(\sqrt{9/2}=3/\sqrt2\)。这也等于把 \((2,1,2)\) 投影到
\((1,1,0)\) 的正交补后所得长度。
\end{solutionBox}

\answerAnchor{MATH1-CALC-0065}
\begin{solutionBox}
\textbf{A}\quad
\[
(1,-1,2)\times(2,0,1)=(-1,3,2).
\]
所得向量与原两向量点积均为 \(0\)。

\textbf{B}\quad
\[
\overrightarrow{AB}=(-1,2,0),\quad
\overrightarrow{AC}=(-1,0,3),\quad
\overrightarrow{AB}\times\overrightarrow{AC}=(6,3,2).
\]
叉积模为 \(7\)，故三角形面积为 \(7/2\)。

\textbf{C}\quad
\[
(1,0,1)\times(0,2,1)=(-2,-1,2),\qquad |(-2,-1,2)|=3.
\]
所有所求单位向量为
\[
\pm\frac13(-2,-1,2).
\]

\textbf{M}\quad
\[
(1,t,1)\times(1,1,-1)=(-t-1,2,1-t),
\]
面积平方为 \(2t^2+6\)。由 \(2t^2+6=6\) 得 \(t=0\)。
\end{solutionBox}

\answerAnchor{MATH1-CALC-0066}
\begin{solutionBox}
\textbf{A}\quad
\[
\det\begin{pmatrix}1&1&0\\0&1&2\\2&0&1\end{pmatrix}=5.
\]

\textbf{B}\quad
\[
V=\frac16\left|
\det\begin{pmatrix}2&0&0\\0&3&0\\1&1&2\end{pmatrix}\right|
=\frac{12}{6}=2.
\]

\textbf{C}\quad
\[
\det\begin{pmatrix}1&0&1\\0&2&1\\1&t&2\end{pmatrix}
=2-t.
\]
共面要求行列式为零，故 \(t=2\)。

\textbf{M}\quad
从 \(A\) 出发，
\[
\overrightarrow{AB}=(-1,1,0),\quad
\overrightarrow{AC}=(-1,0,1),\quad
\overrightarrow{AD}=(0,1,1),
\]
其混合积为 \(2\)，故 \(V=|2|/6=1/3\)。
交换两个顶点只改变有向行列式的符号或等价地改变定向，绝对值不变，
所以几何体积不变。
\end{solutionBox}

\answerAnchor{MATH1-CALC-0067}
\begin{solutionBox}
\textbf{A}\quad
\[
\frac{(1,-2,2)}{\sqrt{1+4+4}}
=\left(\frac13,-\frac23,\frac23\right).
\]

\textbf{B}\quad
\[
\cos^2\gamma=1-\frac14-\frac14=\frac12.
\]
因 \(\gamma\) 为钝角，\(\cos\gamma=-\sqrt2/2\)。

\textbf{C}\quad
比例向量 \((-1,2,2)\) 的模为 \(3\)，故方向余弦为
\[
\left(-\frac13,\frac23,\frac23\right).
\]

\textbf{M}\quad
设 \(\cos\alpha=\cos\beta=c>0\)。由
\[
2c^2+\cos^260^\circ=1
\]
得 \(c=\sqrt6/4\)。故三个方向余弦为
\[
\left(\frac{\sqrt6}{4},\frac{\sqrt6}{4},\frac12\right).
\]
\end{solutionBox}

\answerAnchor{MATH1-CALC-0068}
\begin{solutionBox}
\textbf{A}\quad
\[
(x-2)+2(y+1)-2(z-1)=0,
\]
即 \(x+2y-2z+2=0\)。

\textbf{B}\quad
\[
\overrightarrow{AB}=(1,0,-1),\quad
\overrightarrow{AC}=(0,2,-1),
\]
叉积可取 \((2,1,2)\)。故
\[
2x+y+2(z-1)=0,
\]
即 \(2x+y+2z-2=0\)。三点代回均成立。

\textbf{C}\quad
直线方向为 \((1,1,0)\)，从基点 \((1,0,0)\) 指向 \(Q\) 的向量为
\((-1,0,1)\)，二者叉积可取 \((1,-1,1)\)。故平面为
\[
x-y+z-1=0.
\]

\textbf{M}\quad
所求平面法向须同时垂直于
\((1,1,0)\)、\((0,1,1)\)，故可取
\[
(1,1,0)\times(0,1,1)=(1,-1,1).
\]
又过原点，所求为 \(x-y+z=0\)。
\end{solutionBox}

\answerAnchor{MATH1-CALC-0069}
\begin{solutionBox}
\textbf{A}\quad
\[
(x,y,z)=(-1,2,0)+t(2,-1,3),
\qquad
\frac{x+1}{2}=\frac{y-2}{-1}=\frac z3.
\]

\textbf{B}\quad
\[
x=1,\qquad y=t,\qquad z=-2+2t.
\]
等价地可写 \(x=1,\ y=(z+2)/2\)，没有零分母。

\textbf{C}\quad
两法向叉积为
\[
(1,1,0)\times(0,1,1)=(1,-1,1).
\]
令 \(y=0\) 得线上点 \((1,0,2)\)，故交线为
\[
(x,y,z)=(1,0,2)+t(1,-1,1).
\]

\textbf{M}\quad
两平面交线方向为
\[
(1,1,1)\times(1,-1,2)=(3,-1,-2).
\]
故所求直线为
\[
(x,y,z)=(1,2,0)+t(3,-1,-2).
\]
\end{solutionBox}

\answerAnchor{MATH1-CALC-0070}
\begin{solutionBox}
\textbf{A}\quad
两方向 \((1,2,0)\)、\((2,4,0)\) 平行，但基点差
\((0,0,1)\) 不平行于方向，故为两条不同的平行直线。

\textbf{B}\quad
联立
\[
(t,0,t)=(1,s,1+s)
\]
得 \(t=1,s=0\)，三坐标相容，故两线相交于 \((1,0,1)\)。

\textbf{C}\quad
方向 \((1,0,0)\)、\((0,1,0)\) 不平行；两参数式的 \(z\) 坐标
分别恒为 \(0\)、\(2\)，故无交点，两线异面。

\textbf{M}\quad
\(\vec n_2=2\vec n_1\)，但把 \(\Pi_2\) 除以 \(2\) 后右端为
\(3/2\ne1\)，故 \(\Pi_1\parallel\Pi_2\) 且不重合。
又
\[
(1,2,-1)\cdot(2,-1,0)=0,
\]
故 \(\Pi_1\perp\Pi_3\)。
\end{solutionBox}

\answerAnchor{MATH1-CALC-0071}
\begin{solutionBox}
\textbf{A}\quad
方向 \((1,-1,0)\) 与法向 \((1,1,1)\) 点积为 \(0\)，且基点满足
\(1+1+0=2\)，故直线包含于平面。

\textbf{B}\quad
方向仍与法向点积为 \(0\)，但基点原点不满足平面，故直线严格平行于平面。

\textbf{C}\quad
方向 \(\vec s=(1,2,-1)\)，法向 \(\vec n=(2,-1,2)\)，故
\[
\sin\alpha=\frac{|\vec s\cdot\vec n|}{|\vec s||\vec n|}
=\frac2{3\sqrt6}=\frac{\sqrt6}{9}.
\]
所以 \(\alpha=\arcsin(\sqrt6/9)\)。

\textbf{M}\quad
平面法向为 \((0,0,1)\)，因此
\[
\sin30^\circ=\frac1{\sqrt{k^2+2}}.
\]
于是 \(\sqrt{k^2+2}=2\)，故 \(k=\pm\sqrt2\)。
\end{solutionBox}

\answerAnchor{MATH1-CALC-0072}
\begin{solutionBox}
\textbf{A}\quad
\[
d=\frac{|2-2-2+1|}{\sqrt{1+4+4}}=\frac13.
\]

\textbf{B}\quad
\(z\) 轴方向为 \((0,0,1)\)，点的横向距离为
\[
d=\sqrt{1^2+1^2}=\sqrt2.
\]

\textbf{C}\quad
过 \(P\) 沿法向 \((1,1,1)\) 写
\[
H=(1,0,0)+t(1,1,1).
\]
代入平面得 \(1+3t=0\)，故
\[
H=\left(\frac23,-\frac13,-\frac13\right),\qquad
d=|t|\sqrt3=\frac1{\sqrt3}.
\]

\textbf{M}\quad
两距离分别为 \(|t-3|/\sqrt3\)、\(|t|\)。平方得
\[
(t-3)^2=3t^2
\Longleftrightarrow 2t^2+6t-9=0.
\]
故
\[
t=\frac{-3\pm3\sqrt3}{2}.
\]
两值代回原绝对值等式均成立。
\end{solutionBox}

\answerAnchor{MATH1-CALC-0073}
\begin{solutionBox}
\textbf{A}\quad
第二平面除以 \(2\) 后为
\[
x-2y+2z=-\frac32.
\]
故
\[
d=\frac{|1-(-3/2)|}{\sqrt{1+4+4}}=\frac56.
\]

\textbf{B}\quad
取 \(L_2\) 基点 \(P=(1,2,0)\)，方向
\(\vec s=(1,-1,1)\)，则
\[
\overrightarrow{OP}\times\vec s
=(1,2,0)\times(1,-1,1)=(2,-1,-3).
\]
故两线距离为
\[
d=\frac{\sqrt{2^2+(-1)^2+(-3)^2}}{\sqrt3}
=\sqrt{\frac{14}{3}}.
\]
两方向成比例，而基点差不平行于方向，故两线确为不同的平行直线。

\textbf{C}\quad
\[
\vec s_1\times\vec s_2=(0,0,1),\qquad
\overrightarrow{P_1P_2}=(0,2,3).
\]
故
\[
d=\frac{|(0,2,3)\cdot(0,0,1)|}{1}=3.
\]
混合积非零，确认两线异面。

\textbf{M}\quad
公共法向为 \((0,0,1)\)，故距离为 \(|a|\)。
由 \(|a|=2\) 得 \(a=\pm2\)。此时方向不平行且混合积非零，
两线均为异面直线。
\end{solutionBox}

\answerAnchor{MATH1-CALC-0074}
\begin{solutionBox}
\textbf{A}\quad
\(x^2+y^2+z^2=9\) 是半径 \(3\) 的球面，只含边界；
\(x^2+y^2+z^2\le9\) 是包含内部的闭球体。

\textbf{B}\quad
\(z=2x^2+y^2\ge0\)，故顶点为原点，沿 \(+z\) 开口。
在 \(z=c\) 平面，
\[
2x^2+y^2=c.
\]
\(c>0\) 时为椭圆，\(c=0\) 退化为顶点，\(c<0\) 无实截痕。

\textbf{C}\quad
第一个方程
\(x^2+2y^2+z^2=4\) 表示一个椭球面，是二维曲面；
加入独立平面 \(z=x-y\) 后，表示椭球面与该平面的空间交线，
通常为一条闭合的一维曲线。等式只给边界，不表示椭球体。

\textbf{M}\quad
\(z=4-x^2-y^2\) 是顶点 \((0,0,4)\)、沿 \(-z\) 开口的圆抛物面。
\[
z=c:\quad x^2+y^2=4-c,
\]
仅当 \(c\le4\) 有截痕。对每个 \((x,y)\in\mathbb R^2\) 都有唯一
\(z=4-x^2-y^2\)，故整个曲面在 \(xy\) 平面的投影是全平面。
\end{solutionBox}

\answerAnchor{MATH1-CALC-0075}
\begin{solutionBox}
\textbf{A}\quad
线段上各点到 \(x\) 轴的距离恒为 \(2\)，而轴向范围不变，故
\[
y^2+z^2=4,\qquad -1\le x\le1.
\]
这是半径 \(2\)、长度 \(2\) 的圆柱侧面，不含两个底面。

\textbf{B}\quad
\[
(x-1)^2+\rho^2=4
\Longleftrightarrow
(x-1)^2+y^2+z^2=4,
\]
即以 \((1,0,0)\) 为球心、半径 \(2\) 的球面。上半圆给出非负半径，
旋转一周覆盖整个球面。

\textbf{C}\quad
绕 \(y\) 轴时以
\(\rho=\sqrt{x^2+z^2}\) 替换非负的原 \(x=1+y^2\)，故
\[
x^2+z^2=(1+y^2)^2.
\]
由于 \(1+y^2>0\)，从径向方程平方不会引入负半径分支。

\textbf{M}\quad
\[
\sqrt{y^2+z^2}=\sqrt{x-1}
\Longleftrightarrow
y^2+z^2=x-1,\quad x\ge1.
\]
两边原本都非负，且最终方程自身推出 \(x\ge1\)，所以平方没有引入额外分支。
\end{solutionBox}

\answerAnchor{MATH1-CALC-0076}
\begin{solutionBox}
\textbf{A}\quad
准线是 \(xy\) 平面的椭圆
\[
\frac{x^2}{4}+y^2=1,
\]
半轴为 \(2,1\)；方程缺 \(z\)，母线平行 \(z\) 轴，故为椭圆柱面。

\textbf{B}\quad
方程缺 \(y\)，准线是 \(xz\) 平面的矩形双曲线 \(xz=1\)，
母线平行 \(y\) 轴，故为双曲柱面；准线有两支。

\textbf{C}\quad
方程缺 \(x\)，准线是 \(yz\) 平面的抛物线
\(y=z^2+1\)，母线平行 \(x\) 轴，故为抛物柱面。
准线顶点为 \((y,z)=(1,0)\)；它沿母线方向生成顶线
\(\{(x,1,0):x\in\mathbb R\}\)，曲面沿 \(+y\) 方向开口。

\textbf{M}\quad
这是轴沿 \(y\) 的有限实心圆柱。侧面为
\[
x^2+z^2=4,\quad -1\le y\le2;
\]
两个底面分别为
\[
y=-1,\ x^2+z^2\le4,\qquad
y=2,\ x^2+z^2\le4.
\]
在 \(xz\) 平面的投影是圆盘 \(x^2+z^2\le4\)。
\end{solutionBox}

\answerAnchor{MATH1-CALC-0077}
\begin{solutionBox}
\textbf{A}\quad
配方得
\[
(x-1)^2+4\left(y+\frac12\right)^2+z^2=12.
\]
故为中心 \((1,-1/2,0)\) 的椭球面，沿 \(x,y,z\) 半轴分别为
\[
2\sqrt3,\quad \sqrt3,\quad 2\sqrt3.
\]

\textbf{B}\quad
两正一负且右端为 \(1\)，故为中心 \((1,-1,0)\)、以负号项
\(z\) 为轴的单叶双曲面。固定任意 \(z=c\) 都有椭圆截面。

\textbf{C}\quad
一正两负，故为以 \(z\) 轴为轴的双叶双曲面。
固定 \(z=c\) 时
\[
\frac{x^2}{4}+y^2=\frac{c^2}{9}-1,
\]
仅当 \(|c|\ge3\) 有截面。

\textbf{M}\quad
移轴 \(X=x-1,\ Z=z+2\) 后方程为
\[
X^2+4y^2=Z^2,
\]
故为顶点 \((1,0,-2)\)、轴平行 \(z\) 轴的椭圆锥面（含上下两叶）。
固定 \(z=c\) 时
\[
(x-1)^2+4y^2=(c+2)^2.
\]
所以所有 \(c\in\mathbb R\) 都有截痕：\(c=-2\) 时退化为顶点，
\(c\ne-2\) 时为中心 \((1,0,c)\) 的椭圆。
\end{solutionBox}

\answerAnchor{MATH1-CALC-0078}
\begin{solutionBox}
\textbf{A}\quad
由 \(z=t\) 消元：
\[
x=z^2,\qquad y=z^3,\qquad -1\le z\le2.
\]
范围不可省略，否则会多出原参数区间外的曲线。

\textbf{B}\quad
令
\[
x=2\cos t,\qquad y=\sin t,\qquad 0\le t<2\pi,
\]
则
\[
z=xy=2\cos t\sin t=\sin2t.
\]
半开区间使曲线恰描一周且不重复端点。

\textbf{C}\quad
\[
\vec r(1)=(0,2,1),\qquad
\vec r'(t)=(2t,2,3t^2),
\]
故 \(t=1\) 处切向量为 \((2,2,3)\ne\vec0\)。

\textbf{M}\quad
令
\[
u=\frac{x+y}{\sqrt2},\qquad v=\frac{x-y}{\sqrt2}.
\]
则 \(x^2+y^2=u^2+v^2,\ z=\sqrt2u\)，球面条件化为
\[
3u^2+v^2=2.
\]
取
\[
u=\sqrt{\frac23}\cos t,\qquad v=\sqrt2\sin t,\qquad
0\le t<2\pi,
\]
得到
\[
\vec r(t)=\left(
\frac{\cos t}{\sqrt3}+\sin t,\
\frac{\cos t}{\sqrt3}-\sin t,\
\frac{2\cos t}{\sqrt3}
\right).
\]
代回两方程恒成立；半开区间避免首尾重复。
\end{solutionBox}

\answerAnchor{MATH1-CALC-0079}
\begin{solutionBox}
\textbf{A}\quad
保留 \(x=3\cos t,y=3\sin t\)，消去 \(t\) 得投影曲线
\[
x^2+y^2=9,\qquad z=0.
\]

\textbf{B}\quad
代入 \(z=2x\)：
\[
x^2+y^2=2x
\Longleftrightarrow (x-1)^2+y^2=1.
\]
每个圆周点均可取 \(z=2x\) 回到原曲线，故投影为
\[
(x-1)^2+y^2=1,\qquad z=0.
\]

\textbf{C}\quad
存在实数 \(z\) 的充要条件是上下界有序：
\[
x^2+y^2\le2x
\Longleftrightarrow (x-1)^2+y^2\le1.
\]
因此投影域是以 \((1,0)\) 为圆心、半径 \(1\) 的闭圆盘；
圆盘内每一点都可取 \(z=x^2+y^2\) 回到原立体。

\textbf{M}\quad
由 \(z=y/2\) 得
\[
x=1+\frac{y^2}{4}.
\]
参数范围 \(-1\le z\le1\) 化为 \(-2\le y\le2\)。故完整投影曲线为
\[
x=1+\frac{y^2}{4},\qquad -2\le y\le2,\qquad z=0.
\]
\end{solutionBox}

\answerAnchor{MATH1-CALC-0080}
\begin{solutionBox}
\textbf{A}\quad
（i）只有各变量一次项，直接配方移轴后识别标准曲面，留在本讲；
（ii）要自行消去一般交叉项并求正交变换，应路由到线性代数二次型。

\textbf{B}\quad
\[
4x^2-8x=4(x-1)^2-4,
\]
故
\[
4(x-1)^2+y^2+z^2=5.
\]
这是中心 \((1,0,0)\) 的椭球面，三个半轴为
\(\sqrt5/2,\sqrt5,\sqrt5\)。

\textbf{C}\quad
由 \(x+y=\sqrt2u,\ x-y=\sqrt2v\)，
\[
x^2-y^2=(x+y)(x-y)=2uv.
\]
结果出现 \(uv\) 交叉项，同时含 \(u,v\)。这也提醒：给定正交变换未必
把任意二次式都化为无交叉项标准形。

\textbf{M}\quad
命题错误。正确分流是：高数复习标准二次曲面、配方平移、轴向缩放、
投影和积分应用；含一般交叉项的正交化、特征值与二次型分类放在线性代数
部分复习。“不属于高数常规要求”不等于“数学一整场不考”。
\end{solutionBox}
